\documentclass[12pt]{article}
\usepackage[margin=1in]{geometry}
\usepackage{amsmath,amsthm,amssymb,scrextend}
\usepackage{array}
\setlength{\parindent}{0pt}

\begin{document}
\newcommand{\ds}{\displaystyle}

\textbf{(i)} Both dice first show a six on throw $r$:

$$q^{2r-2}p^2$$

Exactly one die first shows a six on throw $s$, the other first shows a six on throw $r$:

$$q^{2s-2}\cdot 2pq\cdot q^{r-s-1}\cdot p=2p^2q^{r+s-2}$$

\begin{align*}
P(R=r)
&=p^2q^{2r-2}+\sum_{s=1}^{r-1}2p^2q^{r+s-2}\\
&=p^2q^{2r-2}+2p^2q^{r-1}\cdot\frac{1-q^{r-1}}{1-q}\\
&=pq^{r-1}(2-q^{r-1}-q^r)
\end{align*}

\begin{align*}
E(R)
&=\sum_{r=1}^{\infty}r\cdot pq^{r-1}(2-q^{r-1}-q^r)\\
&=p\left(2\sum_{r=1}^{\infty}rq^{r-1}-\sum_{r=1}^{\infty}rq^{2r-2}-\sum_{r=1}^{\infty}rq^{2r-1}\right)\\
&=p\left(\frac{2}{(1-q)^2}-\frac{1}{(1-q^2)^2}-\frac{q}{(1-q^2)^2}\right)\\
&=\frac{3-2p}{p(2-p)}
\end{align*}

\bigskip\hrule\bigskip

\textbf{(ii)}

$$m=\frac{3-2p}{p(2-p)}$$

$$\implies mp^2-(2m+2)p+3=0$$

$$p=\frac{m+1\pm\sqrt{m^2-m+1}}{m}$$

$$\because 0<p\leq1$$

$$\therefore p=\frac{m+1-\sqrt{m^2-m+1}}{m}$$

\bigskip
\hrule
\begin{center}
    This solution was submitted by \textbf{Jon}\\[4pt]
    Published and refined by \textbf{Jon}\\[8pt]
    \textit{Last updated: July 26, 2026}
\end{center}

\end{document}